\documentclass[11pt]{article}
\usepackage[T1]{fontenc}
\usepackage[utf8]{inputenc}
\usepackage{lmodern,microtype,amsmath,amssymb,amsthm,mathtools}
\usepackage[a4paper,margin=28mm]{geometry}
\usepackage{booktabs,array,longtable,enumitem}
\usepackage[authoryear,round]{natbib}
\usepackage[colorlinks=true,allcolors=blue,breaklinks=true]{hyperref}
\usepackage{xurl}
\hypersetup{pdftitle={Grothendieck's questions on tori in apposition, after Kostant and Steinberg},pdfauthor={},pdfsubject={Algebraic groups and representation theory},pdfkeywords={tori in apposition, Borel subgroups, Coxeter elements, Grothendieck, outer automorphisms}}
\setlength{\emergencystretch}{2em}
\setlist{nosep}
\newcommand{\workbenchurl}{https://github.com/KokunoYumeto/algebraic-groups-and-representations-workbench}
\newcommand{\recordurl}{https://github.com/KokunoYumeto/mathematics-workbenches}
\newcommand{\C}{\mathbf C}\newcommand{\Z}{\mathbf Z}\newcommand{\Q}{\mathbf Q}
\newcommand{\g}{\mathfrak g}\newcommand{\tfr}{\mathfrak t}
\newcommand{\Ad}{\operatorname{Ad}}\newcommand{\Aut}{\operatorname{Aut}}
\newcommand{\Out}{\operatorname{Out}}\newcommand{\Inn}{\operatorname{Inn}}
\newcommand{\Lie}{\operatorname{Lie}}\newcommand{\GL}{\operatorname{GL}}
\newcommand{\SL}{\operatorname{SL}}\newcommand{\PGL}{\operatorname{PGL}}
\newcommand{\Sp}{\operatorname{Sp}}\newcommand{\SO}{\operatorname{SO}}
\newcommand{\rank}{\operatorname{rank}}\newcommand{\tr}{\operatorname{Tr}}
\newcommand{\height}{\operatorname{ht}}\newcommand{\diag}{\operatorname{diag}}
\newcommand{\Stab}{\operatorname{Stab}}\newcommand{\id}{\mathrm{id}}
\newtheorem{theorem}{Theorem}[section]
\newtheorem{proposition}[theorem]{Proposition}
\newtheorem{lemma}[theorem]{Lemma}
\newtheorem{corollary}[theorem]{Corollary}
\theoremstyle{definition}\newtheorem{definition}[theorem]{Definition}
\theoremstyle{remark}\newtheorem{remark}[theorem]{Remark}
\title{Grothendieck's questions on tori in apposition,\\ after Kostant and Steinberg}
\author{}
\date{23 November 2026}
\begin{document}
\maketitle
\begin{abstract}
In his letter to Kostant of 22 October 1969, Grothendieck counted pairs of tori in apposition equipped with containing Borel subgroups and asked when their flags are opposite and whether a distinguished class exists. We justify his count, determine exactly the simple types in which every choice is opposite, and prove that a class invariant under all outer automorphisms exists precisely outside the even-rank types $A_{2m}$. The proof uses Kostant's cyclic-element theorems and Steinberg's construction from orthogonal diagram orbits. We specify what naturality under group isomorphisms means: it is equivalent to outer invariance and does not determine a unique class. In type $A_3$, all eight outer-invariant classes are opposite, so opposition still gives no unique choice. We also recall the known answers to the letter's two subregular predictions, with the rank theorem credited to Slodowy through Richardson. The mathematics and exposition were produced by GPT-6 Astra (ChatGPT Pro mode), with same-model mathematical checks; no human has checked them.
\end{abstract}

\section{The letter and the answers}
The first page of Grothendieck's letter asks for a finite classification that rigidifies a simple algebraic group while keeping its outer symmetries visible. The document is cote~123 of the Montpellier archive, scan~22, handwritten folio~21; the letter continues on scan~23, folio~22 \citep{Grothendieck1969}. The numbers of classes and the normalizer extension below are already written there. They are Grothendieck's assertions, for which we give proofs. His questions concern opposition, a distinguished class, and the remaining automorphism groups. The following page suggests two distinctions between simply laced and other root systems at subregular elements; those distinctions fail by a known theorem.

Throughout, $G$ is a connected simple adjoint algebraic group over $\C$. Its Weyl group is $W$, its rank is $r$, its Coxeter number is $h$, and
\[
 z=|Z(G_{\mathrm{sc}})|,
 \qquad U_h=(\Z/h\Z)^\times.
\]
Here $G_{\mathrm{sc}}\to G$ is the simply connected covering. All automorphisms and isomorphisms are over $\C$.

Choose simple roots of a maximal torus $T$ and the coweight $\rho^\vee$ characterized by $\alpha_i(\rho^\vee)=1$. A \emph{principal element} is a conjugate of
\[
 c=\rho^\vee(e^{2\pi i/h}).
\]
An ordered pair $(T_0,T_1)$ is \emph{in principal apposition} if a principal element $c\in T_0$ normalizes $T_1$, has $Z_G(c)^\circ=T_0$, and acts on $\Lie(T_1)$ as a Coxeter element. This is the group formulation of Kostant's apposition construction \citep[Theorems~1.2--1.3, pp.~6; principal-element paragraph, p.~7; Proposition~1.8 and Theorem~1.9, p.~8]{Kostant2010}. The identity component is necessary: a regular semisimple element of an adjoint group can have disconnected centralizer. The order of the two tori matters.

A \emph{quadruple} is
\[
 \mathcal Q=(T_0,B_0,T_1,B_1),\qquad T_i\subset B_i,
\]
where the $B_i$ are Borels and $(T_0,T_1)$ is in apposition. Write $\mathcal X(G)$ for the set of their simultaneous inner-conjugacy classes. Borels are \emph{opposite} if their intersection is a maximal torus.

\begin{theorem}\label{thm:main}
For every such $G$ the following statements hold.
\begin{enumerate}[label=\textup{(\roman*)}]
\item All ordered apposition pairs are conjugate. Their simultaneous normalizer has order $hz\varphi(h)$ and acts freely on the containing Borel pairs. Thus
\[
 |\mathcal X(G)|=\frac{|W|^2}{hz\varphi(h)}.
\]
If a generating principal element of $T_0\cap N_G(T_1)$ is included, the number is $|W|^2/(hz)$.
\item Given either Borel choice, an opposite choice for the other torus exists. Every Borel choice is opposite if and only if the type is $A_{p-1}$ for a prime $p$, $B_2=C_2$, or $G_2$.
\item Every quadruple has trivial inner stabilizer, and
\[
 \Aut(G;\mathcal Q)\simeq\Stab_{\Out(G)}[\mathcal Q].
\]
There is a class fixed by all of $\Out(G)$ if and only if the type is not $A_{2m}$, $m\geq1$.
\end{enumerate}
\end{theorem}

The theorem makes the letter's finite rigidification precise. The positive assertion in (ii) holds for arbitrary two maximal tori, without apposition. Its universal assertion is more restrictive: a cyclic relation between the tori need not make all of their flags transverse. The prime-degree case is exactly the nonvanishing theorem for Fourier minors attributed to Chebotar\"ev; we use the published formulation in \citet[Lemma~1.3, p.~123]{Tao2005}.

Naturality requires a separate statement. A rule assigning $x_G\in\mathcal X(G)$ is \emph{natural under isomorphisms} when every isomorphism $f:G\to H$ carries $x_G$ to $x_H$. This definition concerns classes, not a representative fixed as four subgroups.

\begin{theorem}\label{thm:natural}
For a fixed simple Dynkin type, evaluation on one group identifies natural rules under isomorphisms with $\mathcal X(G)^{\Out(G)}$. Consequently such rules exist exactly outside $A_{2m}$. They need not be unique. For $G=\PGL_4$ there are exactly eight such choices, all with opposite Borels. Compatibility with central isogenies does not reduce that ambiguity. The standard symplectic folding selects exactly four of these eight classes. Requiring all four subgroups to be defined over the split real form cannot give a rule in type $A_{n-1}$ for any $n\geq3$.
\end{theorem}

The theorem answers a specified equivariance question. It does not identify a geometrically preferred one of the remaining choices, or a choice compatible with all diagram foldings; those are different requirements, formulated in Section~\ref{sec:natural}. In particular an arbitrary ordering of a finite classification would not establish geometric canonicity.

\paragraph{Credit and context.}
Kostant's principal and cyclic elements carry the reduction to a finite problem. His original Coxeter statement is Theorem~8.6 and its corollary \citep[p.~1026]{Kostant1959}; the directly used modern account also states the cyclotomic normalizer quotient in equation~(0.6) \citep[p.~3]{Kostant2010}. Steinberg's proof of Theorem~9.4 \citep[pp.~72--73]{Steinberg1965} supplies the orthogonal diagram-orbit construction and the same $A_{2m}$ exception, there for rational cross-sections. We apply that construction to an invariant quadruple and prove the converse here. The remaining opposition cases are established by explicit classical matrices and exceptional weight calculations.

Apposition also enters the Stokes geometry of the $tt^*$-Toda equations: \citet[\S3, pp.~9--10; \S5]{GuestHo2019} use Kostant's Cartan and Coxeter structures to describe Stokes data. \citet[Introduction and \S6]{Elashvili2013} extend the study of cyclic elements beyond principal nilpotents. These results explain why cyclic Cartans continue to be useful, but their stated conclusions are different from the finite flag classification considered here.

The recent machine-assisted archive survey of \citet[Part~II, Problem~7, pp.~391--392]{Hua2026} restates the opposition and distinguished-class questions and gives an argument for the count. It explicitly leaves the general relative-position question open; the formal checks reported there concern numerical integrality of the count. We give the geometric arguments and the complete opposition classification here. No claim of priority is made for the deductions in this paper.

\section{Proof overview}
The main difficulty is that two different symmetries are present. A principal element gives a Coxeter action on the second torus, while the Borels are chambers for two distinct root systems. First we determine the common normalizer, including its finite torus kernel. Kostant's conjugacy theorem gives every cyclotomic multiplier. The integral identity $(1-C)P=Q$ identifies the remaining kernel with the center of the simply connected group. Freeness then follows from the trivial intersection of the two tori; dividing by the normalizer order is justified only after this step.

Opposition is tested in representations. In type $A$ the change of eigenbasis is a Fourier matrix. The symplectic and orthogonal models keep track of which Fourier flags are allowed by the form. For $E_6$ and $E_7$, a short Coxeter orbit of an extremal weight forces a zero matrix coefficient. For $F_4$ and $E_8$, a highest-weight line in $\bigwedge^2\g$ gives a vanishing determinant. The positive $G_2$ case requires both fundamental representations.

For outer symmetry, commuting rank-one subgroups attached to diagram orbits give a fixed Coxeter representative outside $A_{2m}$. A torus in the connected fixed group then provides a fixed regular cocharacter, and hence a stable Borel. The exclusion in even-rank type $A$ is an elementary obstruction to two simultaneous affine reversals. Counting those reversals also shows why invariance and opposition leave several possible distinguished classes. Adjointness is used in the torus intersection and in the order of a principal element; characteristic zero supplies the cyclic-element and representation theorems.

\section{The common normalizer and Grothendieck's count}\label{sec:count}
This section proves the count together with the finite group that makes it meaningful. We use the following specific consequences of Kostant's theorems. If $\theta=\sum n_i\alpha_i$ is the highest root and $\zeta=e^{2\pi i/h}$, then
\begin{equation}\label{eq:cyclic}
 x=\sum_{i=1}^r b_i e_{\alpha_i}+b_0e_{-\theta},\qquad b_i\ne0,
 \quad c=\rho^\vee(\zeta)
\end{equation}
is regular semisimple. Its torus $T_1=Z_G(x)$ is Coxeter for $c$, and $\C x$ is the one-dimensional $\zeta$-eigenspace in $\Lie(T_1)$. A homogeneous invariant of degree $h$ is nonzero on $x$. Principal elements form one conjugacy class, every Coxeter representative is principal, and $T_0=Z_G(c)^\circ$ is transitive on the Cartans Coxeter for $c$ \citep[Proposition~1.1, Theorems~1.2--1.4, and Theorem~1.9, pp.~5--8]{Kostant2010}.

\begin{lemma}\label{lem:rigidity}
All ordered apposition pairs are conjugate, and for the pair in \eqref{eq:cyclic},
\[
 T_0\cap N_G(T_1)=\langle c\rangle,\qquad T_0\cap T_1=1.
\]
Every quadruple has trivial inner stabilizer.
\end{lemma}
\begin{proof}
Conjugate a principal element of the given pair to $c$. The first torus becomes $T_0$, and Kostant's torus transitivity carries the second to $T_1$.

If $t\in T_0\cap N_G(T_1)$, its commutation with $c$ implies $\Ad(t)x=\lambda x$. Comparing the nonzero coefficients in \eqref{eq:cyclic} gives $\alpha_i(t)=\lambda$ and $(-\theta)(t)=\lambda$. As $\sum n_i=h-1$, we obtain $\lambda^h=1$. In an adjoint group the simple roots form a basis of $X^*(T_0)$; thus $t=c^j$ for the unique $j$ such that $\lambda=\zeta^j$. Conversely $c$ normalizes $T_1$. An element also in $T_1$ fixes $x$, so has $\lambda=1$ and equals $1$.

An element normalizing $T_i$ and preserving $B_i\supset T_i$ lies in $T_i$, since the Weyl action on its chambers is free. A stabilizer of the quadruple therefore lies in $T_0\cap T_1$.
\end{proof}

Set $W_i=N_G(T_i)/T_i$ and let $C\in W_1$ be the image of $c$. Define
\begin{align*}
 A_h&=\{w\in W_0:w(c)=c^a\text{ for some }a\in U_h\},
 &\mu_0(w)&=a,\\
 N_h&=N_{W_1}(\langle C\rangle),
 &wCw^{-1}&=C^{\mu_1(w)}.
\end{align*}
These definitions are finite root-datum tests. For example, with $P^\vee=X_*(T_0)$,
\[
 A_h=\{w:w\rho^\vee-a\rho^\vee\in hP^\vee\text{ for some }a\in U_h\}.
\]

\begin{proposition}\label{prop:normalizer}
Both multiplier maps are surjective and the two Weyl projections identify
\[
 K:=N_G(T_0)\cap N_G(T_1)\simeq A_h\times_{U_h}N_h.
\]
There is an exact sequence
\begin{equation}\label{eq:extension}
 1\longrightarrow\langle c\rangle\times T_1^c
 \longrightarrow K\longrightarrow U_h\longrightarrow1,
 \qquad T_1^c\simeq Z(G_{\mathrm{sc}}).
\end{equation}
The last isomorphism is one of finite diagonalizable groups; no preferred isomorphism or splitting of the extension is asserted.
\end{proposition}
\begin{proof}
First $Z_{W_1}(C)=\langle C\rangle$. A commuting Weyl element preserves the line $\C x$, so acts on it by some $\lambda$. The nonzero invariant of degree $h$ gives $\lambda^h=1$. This scalar determines the element, since a Weyl element fixing a regular complex vector is trivial. Indeed write that vector as $a+ib$ in the real reflection representation and move $a$ into a closed chamber. Its stabilizer is the reflection group of walls through $a$; regularity makes $b$ regular for this smaller root system, so its stabilizer there is trivial. Since $C$ realizes every $h$th-root scalar, the centralizer is exactly its cyclic group.

Every primitive eigenvalue of $C$ has a regular eigenvector. To see this without an additional classification, choose a nonzero $\zeta$-eigenvector over $\Q(\zeta)$ in the rational cocharacter representation. Its root values are nonzero because its line is $\C x$. Galois conjugation gives the same assertion for $\zeta^a$, $a\in U_h$.

The group $K$ normalizes the intrinsic group $T_0\cap N_G(T_1)=\langle c\rangle$, giving a multiplier map $\mu:K\to U_h$. Fix $a\in U_h$. The element $c^a$ is regular of order $h$, hence principal. Choose $n$ with $ncn^{-1}=c^a$. It normalizes $T_0$. On $nT_1n^{-1}$, $c$ acts as $C^{a^{-1}}$, which has a regular $\zeta$-eigenvector by the preceding paragraph. This vector lies in the $\zeta$-eigenspace of $\Ad(c)$ in $\g$, so Kostant's cyclic theorem identifies its Cartan as one of the Cartans Coxeter for $c$. Torus transitivity supplies $t\in T_0$ such that $tn\in K$. Thus $\mu$ is onto.

Let $F=T_1^c$. It centralizes $c$, so normalizes $Z_G(c)^\circ=T_0$ and belongs to $K$. If $k\in\ker\mu$, its image in $W_1$ centralizes $C$, hence is $C^j$. Therefore $c^{-j}k\in T_1^c$. Conversely $\langle c\rangle$ and $F$ centralize each other and lie in the kernel. Their intersection is trivial by Lemma~\ref{lem:rigidity}. This proves the kernel description.

For its order, let $Q\subset P$ be the root and weight lattices for $T_1$, and let $D=s_1\cdots s_r$ be the Coxeter action on characters. For the fundamental weights $\omega_i$,
\[
 (1-D)\omega_i=s_1\cdots s_{i-1}\alpha_i.
\]
The right-hand vectors form a triangular integral basis of $Q$. Thus $(1-D)P=Q$, and $1-D$ induces
\[
 P/Q\simeq Q/(1-D)Q=X^*(F).
\]
Since $P/Q$ is the character group of $Z(G_{\mathrm{sc}})$, we obtain the final assertion of \eqref{eq:extension} and $|K|=hz\varphi(h)$.

Finally $K\to W_0\times W_1$ is injective by the trivial torus intersection. The multipliers of its two projections agree. Every $w_0\in A_h$ lifts to $K$ by the same torus adjustment used above. Given an element of the indicated fiber product, lift its first component. The discrepancy in the second component centralizes $C$, so multiplication by a power of $c$ corrects it without changing the first component. This proves the fiber-product assertion and the surjectivity of both multipliers.
\end{proof}

After choosing one Borel through each torus, the Borel pairs form the $K$-set $W_0\times W_1$, with
\[
 k\cdot(w_0,w_1)=(p_0(k)w_0,p_1(k)w_1).
\]
Lemma~\ref{lem:rigidity} makes the action free. Its quotient is exactly $\mathcal X(G)$, proving Theorem~\ref{thm:main}(i). Including a generator of $\langle c\rangle$ leaves, after normalizing it to $c$, the group $\ker\mu$, giving the second count. Also the map
\[
 \Aut(G;\mathcal Q)\longrightarrow\Out(G)
\]
is injective. An outer class preserving $[\mathcal Q]$ has a representative which preserves $\mathcal Q$ after composing with an inner automorphism, so its image is exactly the asserted stabilizer.

\begin{center}
\begin{tabular}{lrr@{\qquad}lrr}
\toprule
Type&$h$&$z$&Type&$h$&$z$\\\midrule
$A_r$&$r+1$&$r+1$&$E_6$&12&3\\
$B_r,C_r$&$2r$&2&$E_7$&18&2\\
$D_r$&$2r-2$&4&$E_8$&30&1\\
$G_2$&6&1&$F_4$&12&1\\\bottomrule
\end{tabular}
\end{center}
For example, the numbers of classes are $2$ in $A_2$, $18$ in $A_3$, $4$ in $B_2$, $12$ in $G_2$, and $768$ in $D_4$.

\section{Opposition in the classical groups}\label{sec:classical}
We first separate existence of an opposite pair from opposition of every pair.
\begin{lemma}\label{lem:some}
If $G$ is connected reductive over $\C$, $B_0$ is a Borel, and $S$ is a maximal torus, some Borel containing $S$ is opposite to $B_0$.
\end{lemma}
\begin{proof}
Choose opposite $B,B^-$ through $S$, write $B_0=gBg^{-1}$, and use the refined Bruhat decomposition
\[
 g=u\dot w b,\qquad u\in U\cap\dot wU^-\dot w^{-1},\quad b\in B
\]
\citep[\S2.4, p.~52]{Steinberg1965}. Put $D=\dot wB^-\dot w^{-1}$. It contains $S$ and $u$. Conjugation by $u$ carries the opposite pair $(\dot wB\dot w^{-1},D)$ to $(B_0,D)$.
\end{proof}

\subsection{Type \texorpdfstring{$A$}{A} and Fourier minors}
Put $n=r+1$, $\zeta=e^{2\pi i/n}$, and let $F=(\zeta^{ja})_{j,a\in\Z/n}$. In $\PGL_n$ take $T_0$ diagonal and $T_1=FT_0F^{-1}$. The diagonal element $\diag(1,\zeta,\ldots,\zeta^{n-1})$ cyclically permutes the Fourier eigenlines, giving apposition. This also follows by taking the cyclic shift as $x$ in \eqref{eq:cyclic}.

Borels through the two tori correspond to orderings $\pi,\sigma$ of $\Z/n$. Their flags are opposite precisely when, for $1\leq k<n$,
\begin{equation}\label{eq:minor}
 \det F_{\{\pi(k),\ldots,\pi(n-1)\},\{\sigma(0),\ldots,\sigma(n-k-1)\}}\ne0.
\end{equation}
Indeed this determinant tests whether the initial Fourier subspace of dimension $n-k$ maps isomorphically to the quotient by the initial coordinate subspace of dimension $k$.

For prime $n$, every square minor is nonzero by Chebotar\"ev's theorem \citep[Lemma~1.3, p.~123; author-version proof, pp.~3--4]{Tao2005}. For composite $n=ab$ with $a,b>1$, the rows $0,a$ and columns $0,b$ give an all-ones $2\times2$ submatrix. Put them respectively last and first in the orderings. Equation~\eqref{eq:minor} fails. This proves the complete type-$A$ assertion.

\subsection{Type \texorpdfstring{$C$}{C}}
For $n=2r$, use the alternating form
\[
 \langle e_i,e_j\rangle=(-1)^i\mathbf1_{i+j=n-1}
\]
and the cyclic shift $Xe_j=e_{j+1}$, indices modulo $n$. It is skew-adjoint. For simple roots $\alpha_i=\varepsilon_i-\varepsilon_{i-1}$ ($1\leq i<r$) and $\alpha_r=-2\varepsilon_{r-1}$, its entries are the nonzero simple-root vectors
\[
 E_{i,i-1}+E_{n-i,n-i-1},\quad E_{r,r-1},
\]
and the negative-highest-root vector $E_{0,n-1}$. Thus its diagonal and centralizing tori give apposition. Its eigenvectors $v_a=\sum_j\zeta^{ja}e_j$ satisfy
\[
 Xv_a=\zeta^{-a}v_a,\qquad
 \langle v_a,v_b\rangle=n\zeta^{-b}\mathbf1_{b=a+r}.
\]
The coordinate dual pairs are $i,n-1-i$; the Fourier dual pairs are $a,a+r$.

A Borel flag chooses one member of each dual pair, in any order, then appends the reversed partners. These are exactly the complete selfdual flags. Two such flags are opposite in the symplectic group exactly when they are opposite as linear flags. For the nontrivial direction, opposite complete linear flags give lines $L_i=F_i\cap G_{n+1-i}$. Selfduality implies $\langle L_i,L_j\rangle=0$ unless $i+j=n+1$. Nondegeneracy makes the remaining pairings nonzero, and the common symplectic stabilizer is the diagonal maximal torus in these lines.

For $r=2$, the one-entry minors of $F_4$ and its inverse are nonzero, so the minors of orders $1$ and $3$ are nonzero. In a relevant $2\times2$ minor the two Fourier columns are not dual, hence have odd difference. Its determinant is a nonzero scalar times $1-\zeta^{(v-u)(b-a)}$, where $u\ne v$ are its rows and $b-a$ is odd. It cannot vanish modulo $4$. All choices are therefore opposite. The rank-one case is equally immediate.

For $r\geq3$, begin the coordinate flag with $n-1,r-1$ and the Fourier flag with $0,2$. These are nondual pairs, so both flags extend. The last two coordinate rows are $r,0$, and their minor in columns $0,2$ is all ones. The flags are not opposite.

\subsection{Types \texorpdfstring{$B$ and $D$}{B and D}}
For $B_r$, $r\geq3$, put $m=2r$ and use $(e_i,e_j)=(-1)^i\mathbf1_{i+j=m}$ on $\C^{m+1}$. The skew-adjoint matrix
\[
 X=\sum_{i=0}^{m-1}E_{i+1,i}+E_{0,m-1}+E_{1,m}
\]
is cyclic: relative to $\alpha_i=\varepsilon_i-\varepsilon_{i-1}$ and $\alpha_r=-\varepsilon_{r-1}$, the subdiagonal entries are the simple-root vectors, while the last two entries have root $-\theta=\varepsilon_0+\varepsilon_1$. Choose $a^m=2$ and $\lambda_j=a\zeta^j$. An eigenbasis is
\[
 v_j=e_0+\sum_{i=1}^{m-1}\lambda_j^{m-i}e_i+e_m,
 \quad Xv_j=\lambda_jv_j,
 \qquad w=e_0-e_m,\quad Xw=0.
\]
Direct substitution gives $(v_j,v_b)=2m\mathbf1_{b=j+r}$, $(w,v_j)=0$, and $(w,w)=-2$. Hence the Borel flags are isotropic flags in the indicated dual pairs, with the zero-weight line in the middle. Opposition implies opposition of their complete linear flags, by conjugation to upper and lower orthogonal triangular flags.

Begin the coordinate isotropic flag with $e_{m-1},e_{r-1}$ and the second with $v_0,v_2$. The quotient by the perpendicular of the first two coordinate lines is measured in rows $1,r+1$. Its matrix on $v_0,v_2$ is
\[
 \begin{pmatrix}a^{m-1}&a^{m-1}\zeta^{-2}\\
 a^{r-1}&a^{r-1}\zeta^{-2}\end{pmatrix},
\]
which is singular. This gives nonopposition. The small cases follow from $B_1=A_1$ and $B_2=C_2$.

For $D_r$, $r\geq4$, work in $\SO_{2r}$ and then divide by its center. Set $n=2r$ and $(e_i,e_{n-1-i})=s_i$, with $s_i=(-1)^i$ for $i<r$ and $s_{n-1-i}=s_i$. With the roots $\alpha_i=\varepsilon_i-\varepsilon_{i-1}$ and $\alpha_r=-\varepsilon_{r-2}-\varepsilon_{r-1}$, take
\[
 X=\sum_{i=1}^{r-1}(E_{i,i-1}+E_{n-i,n-i-1})
       +E_{r,r-2}+E_{r+1,r-1}+E_{0,n-2}+E_{1,n-1}.
\]
These are skew-adjoint simple-root and negative-highest-root vectors, so $X$ is cyclic. Its kernel is the nondegenerate plane spanned by
\[
 u=e_0-e_{n-1},\qquad v=e_{r-1}-e_r,
\]
with $(u,u)=-2$, $(v,v)=-2(-1)^{r-1}$ and $(u,v)=0$. If $\lambda^2=(-1)^r$, its isotropic lines are $w_\pm=u\pm\lambda v$. The centralizing torus preserves them: it is connected, and the orthogonal torus of this plane, acting trivially on the complement, centralizes $X$.

The coordinate eigenline $\C e_1$ is orthogonal to $\C w_+$. Each can be the first line of a torus-containing isotropic Borel flag. First lines for opposite orthogonal Borels pair nontrivially, as is seen on standard opposite flags. The resulting Borels are not opposite. Their images in the adjoint group give the required example.

\section{Opposition in the exceptional groups}\label{sec:exceptional}
We now give short exact certificates for all exceptional types. Representations may be taken on $G_{\mathrm{sc}}$: their projective highest lines and Borel subgroups descend through the central isogeny. We use the highest-weight theorem and complete reducibility over $\C$ \citep[Proposition~22.17, Theorems~22.18--22.20, pp.~467--468; Proposition~22.41, p.~477]{Milne2017}.

\subsection{Short extremal orbits: \texorpdfstring{$E_6$ and $E_7$}{E6 and E7}}
\begin{lemma}\label{lem:short}
Suppose a Coxeter element $C$ has an orbit of length $d<h$ on the extremal weights $W\omega_i$ of a fundamental representation. Then some containing Borel pair for an apposition pair is not opposite.
\end{lemma}
\begin{proof}
Choose the Coxeter action on $T_1$, lift it to a principal element $c$, and put $T_0=Z_G(c)^\circ$. Choose a positive system for $T_0$ with $\alpha_j(c)=\zeta$. The representation has a fundamental highest weight $\lambda$ for this base, since inner conjugacy transports based tori.

Let $u$ span a $T_1$ extremal weight line fixed by $C^d$. For a lift $\widetilde c$ to the simply connected group, $\widetilde c^d u=bu$. On the two $T_0$ extremal lines of weights $\lambda$ and $s_j\lambda=\lambda-\alpha_j$, where $\lambda=\omega_j$, the eigenvalues of $\widetilde c^d$ have ratio $\zeta^d\ne1$. At least one, say $a$ on weight $\nu$, differs from $b$. A nonzero dual weight vector $f$ supported on $V_\nu$ satisfies
\[
 a f(u)=f(\widetilde c^d u)=b f(u),\qquad f(u)=0.
\]
Choose Borels through $T_0,T_1$ making $f,u$ highest vectors of $V^*,V$. If the Borels were opposite, these lines would pair nontrivially: for standard opposite Borels they are opposite extremal dual lines, and conjugation preserves evaluation. This is a contradiction.
\end{proof}

Here are integral certificates for the required orbits. In fundamental-weight coordinates use the simply laced Cartan matrix with edges
\[
 (1,3),(2,4),(3,4),(4,5),(5,6),
 \quad\text{and }(6,7)\text{ for }E_7.
\]
A reflection has formula $s_j(v)_k=v_k-v_jA_{jk}$; let $C=s_1\cdots s_r$, with the rightmost factor applied first.

For $E_6$, the reflection word $1,3,4,2,5,4$ in application order sends $\omega_1$ to $\mu=(0,0,1,-1,0,1)$. Its cycle is
\begin{align*}
 (0,0,1,-1,0,1)&\longmapsto(-1,0,0,1,-1,0)\\
 &\longmapsto(1,0,-1,0,1,-1)
 \longmapsto(0,0,1,-1,0,1).
\end{align*}
Thus $d=3<h=12$. For $E_7$, the word $7,6,5,4,2,3,1,4,3,5,4,6$ sends $\omega_7$ to $\mu=(0,1,0,-1,1,-1,1)$, and
\[
 C\mu=(0,-1,0,1,-1,1,-1),\qquad C^2\mu=\mu.
\]
Here $d=2<h=18$. Every equality follows by the displayed integer reflection formula. Lemma~\ref{lem:short} proves nonopposition in both types.

\subsection{Highest wedges: \texorpdfstring{$F_4$ and $E_8$}{F4 and E8}}
Let $\kappa$ be the Killing form and let its induced form on $\bigwedge^2\g$ be
\[
 \kappa_2(a\wedge b,u\wedge v)
 =\kappa(a,u)\kappa(b,v)-\kappa(a,v)\kappa(b,u).
\]
Suppose distinct roots $\theta,\delta$ have the following properties: $\theta+\alpha_i$ is never a root; the only root among $\delta+\alpha_i$ is $\theta$; and $-1\in W$. Then $e_\theta\wedge e_\delta$ is killed by every positive simple-root operator. Complete reducibility shows that it generates an irreducible highest-weight submodule $V$. Its pairing with $e_{-\theta}\wedge e_{-\delta}$ is nonzero, so $\kappa_2|_V$ is nondegenerate. Consequently the highest lines for opposite Borels in $V$ always pair nontrivially.

Let $\alpha,\beta$ for $T_0$ be an unordered Weyl translate of $\theta,\delta$. Suppose the $T_1$ roots $\mu,C^d\mu$ form another such translate and
\begin{equation}\label{eq:wedgecondition}
 h\mid d\bigl(\height(\alpha)-\height(\beta)\bigr).
\end{equation}
Choose $u\in\g_\mu(T_1)$ nonzero. The two highest-line vectors
\[
 z_0=e_\alpha\wedge e_\beta,\qquad
 z_1=u\wedge\Ad(c^d)u
\]
lie in the same module $V$. If $a=\kappa(e_\alpha,u)$ and $b=\kappa(e_\beta,u)$, invariance gives
\[
 \kappa_2(z_0,z_1)=ab\left(
 \zeta^{-d\height(\beta)}-\zeta^{-d\height(\alpha)}\right)=0.
\]
Their Borels cannot be opposite. Choosing the Coxeter word on $T_1$ first and conjugating its principal lift to $\rho^\vee(\zeta)$ justifies these simultaneous coordinates on the two tori.

For $F_4$, use the roots $\pm e_i$, $\pm e_i\pm e_j$, and $\tfrac12(\pm e_1\pm e_2\pm e_3\pm e_4)$, with simple roots
\[
 a_1=e_2-e_3,\quad a_2=e_3-e_4,\quad a_3=e_4,\quad
 a_4=\tfrac12(e_1-e_2-e_3-e_4).
\]
Take $\theta=e_1+e_2$ and $\delta=e_1+e_3$. The raising conditions follow from this root list; $-1$ is the product of reflections in the four short roots $e_i$. The reflections $s_2,s_3,s_4$ fix $\theta$ and send $\delta$ successively to $e_1+e_4$, $e_1-e_4$, and $e_2+e_3$. Thus take $\alpha=e_2+e_3$, $\beta=\theta$. Their heights are $5,11$, since the height functional is $(8,3,2,1)$. For $C=s_1s_2s_3s_4$,
\[
 C\theta=e_1+e_3,\qquad C^2\theta=e_1+e_4.
\]
The pair $\theta,C^2\theta$ is the $s_2$-translate of $\theta,\delta$. With $h=12$, $d=2$, the divisibility in \eqref{eq:wedgecondition} holds.

For $E_8$, use the roots $\pm e_i\pm e_j$ and the half-vectors $\tfrac12\sum\varepsilon_ie_i$ with an even number of minus signs. Set
\begin{align*}
 a_1&=\tfrac12(e_1-e_2-e_3-e_4-e_5-e_6-e_7+e_8),\\
 a_2&=e_1+e_2,\quad a_3=e_2-e_1,\quad
 a_i=e_{i-1}-e_{i-2}\quad(4\leq i\leq8).
\end{align*}
Take $\theta=e_7+e_8$, $\delta=e_6+e_8$. The root list verifies the raising conditions. The reflections in $e_1-e_2,e_1+e_2,\ldots,e_7-e_8,e_7+e_8$ multiply to $-1$. The word $s_7,s_6,s_5,s_4,s_2$ fixes $\theta$ and moves $\delta$ through
\[
 e_5+e_8, e_4+e_8, e_3+e_8, e_2+e_8, e_8-e_1.
\]
Take $\alpha=e_8-e_1$, $\beta=\theta$. Their heights are $23,29$, as the height functional is $(0,1,2,3,4,5,6,23)$. For $C=s_1\cdots s_8$ and $\mu=e_2+e_8$, direct reflection gives
\[
 C^j\mu=e_{j+2}+e_8\quad(0\leq j\leq5).
\]
The pair $\mu,C^5\mu$ is a Weyl translate of $\theta,\delta$ by the first four reflections in the preceding word. Now $h=30$ and $5(29-23)=30$. This proves nonopposition in $E_8$ as well. These explicit Euclidean root systems agree with \citet[\S24(l), pp.~539--542]{Milne2017}.

\subsection{The positive case \texorpdfstring{$G_2$}{G2}}
Let $\alpha,\beta$ be short and long simple roots, with squared lengths $2,6$ and inner product $-3$. The fundamental highest weights are $\omega_1=2\alpha+\beta$ and $\omega_2=3\alpha+2\beta$. The Weyl dimension formula gives
\[
 \dim L(\omega_1)=\frac21\frac33\frac54\frac75\frac96\frac{12}9=7
\]
\citep[Proposition~28.1, p.~125]{Etingof}. Its weights are the six short roots and zero, each once. Indeed the extremal orbit gives six multiplicity-one weights, and the remaining weight must be Weyl-fixed. Since $w_0=-1$, this representation is self-dual. Its invariant form $B$ is symmetric because its dimension is odd.

A principal $\mathfrak{sl}_2$ has weights $-6,-4,-2,0,2,4,6$ on it, so the representation restricts to the single irreducible module of highest weight six \citep[\S33.3, p.~149]{Etingof}. Choosing weight vectors gives
\[
 N=e_\alpha+e_\beta=\sum_{i=0}^5E_{i+1,i},\qquad
 B(e_i,e_j)=(-1)^i\delta_{i+j,6},
\]
with weights, in order,
\[
 -(2\alpha+\beta),\ -\alpha-\beta,\ -\alpha,\ 0,\
 \alpha,\ \alpha+\beta,\ 2\alpha+\beta.
\]
The grade $-5$ subspace of $\mathfrak{so}(V,B)$ is the single line spanned by $Y=E_{0,5}+E_{1,6}$. The negative-highest-root space of $\g_2$ is nonzero and has this grade, so it is that line. Therefore $X=N+Y$ is a cyclic element of $\g_2$.

For $a^6=2$, $\zeta=e^{2\pi i/6}$, $\lambda_j=a\zeta^j$, its eigenvectors are
\[
 v_j=e_0+\sum_{i=1}^5\lambda_j^{6-i}e_i+e_6,
 \quad Xv_j=\lambda_jv_j,
 \qquad w=e_0-e_6,\quad Xw=0.
\]
The principal element rotates the six nonzero eigenlines in a Coxeter cycle. A Coxeter element rotates the short-root hexagon by $60$ degrees; hence consecutive $v_j,v_{j+1}$ have weights summing to a long root.

The adjoint representation is a submodule of $\bigwedge^2V$, via $\g_2\subset\mathfrak{so}(V,B)$. Its six extremal long-root lines for the first torus are
\[
 \C(e_i\wedge e_k),\qquad
 (i,k)\in\{(0,1),(0,2),(1,4),(2,5),(4,6),(5,6)\}.
\]
They are obtained by adding the displayed weights. Each has multiplicity one in $\bigwedge^2V$. For the second torus they are $\C(v_j\wedge v_{j+1})$, $j\in\Z/6$. The induced form $B_2$ is nondegenerate on this adjoint submodule because it pairs the first and last coordinate wedge lines nontrivially.

All pairings $B(e_i,v_j)$ for nonzero coordinate weights are nonzero. For the adjoint lines put $(q_0,\ldots,q_6)=(0,5,4,3,2,1,0)$. Up to a nonzero sign their pairings are
\[
 \det\begin{pmatrix}
 \lambda_j^{q_{6-i}}&\lambda_{j+1}^{q_{6-i}}\\
 \lambda_j^{q_{6-k}}&\lambda_{j+1}^{q_{6-k}}
 \end{pmatrix}
 =\lambda_j^{q_{6-i}+q_{6-k}}
   (\zeta^{q_{6-k}}-\zeta^{q_{6-i}}).
\]
The only repeated residues among the $q_i$ are at $0,6$, and none of the listed pairs is $(0,6)$. Thus all these pairings are nonzero too.

Finally, these two tests suffice. For a fixed positive Borel and its two fundamental highest vectors $z_i$, write $g=b_1\dot u b_2$ in Bruhat form. Invariance of the forms and stability of the highest lines imply
\[
 (z_i,gz_i)\ne0\quad\Longleftrightarrow\quad u\omega_i=-\omega_i.
\]
The two fundamental weights are a basis, so both nonvanishings force $u=-1=w_0$, exactly the opposite Bruhat cell. Every containing Borel pair passes these tests. This proves the $G_2$ assertion and completes Theorem~\ref{thm:main}(ii).

\section{Outer automorphisms and the affine obstruction}\label{sec:outer}
This section proves existence and nonexistence of a fully outer-invariant class. A pinning identifies $\Gamma=\Out(G)$ with the automorphism group of the Dynkin diagram and realizes it as automorphisms of $G$ preserving the pinning \citep[Propositions~23.44--23.45 and Corollaries~23.46--23.47, pp.~497--498]{Milne2017}.

\subsection{The construction after Steinberg}
Outside $A_{2m}$, every $\Gamma$-orbit of simple roots consists of mutually orthogonal roots. This follows directly from the diagrams, including the three outer vertices of $D_4$. It is the orthogonality used by \citet[proof of Theorem~9.4, pp.~72--73]{Steinberg1965}.

Let $\phi_i:\SL_2\to G$ be the pinned rank-one homomorphisms. For a diagram orbit $O$, their images commute: the positive and negative root subgroups commute by the root commutator relations, and each root is trivial on the other coroot torus. Thus
\[
 \phi_O(A)=\prod_{i\in O}\phi_i(A)
\]
is a homomorphism into $H=(G^\Gamma)^\circ$, fixed pointwise by $\Gamma$. Choose any ordering of the orbits and put
\[
 n=\prod_O\phi_O\begin{pmatrix}0&1\\-1&0\end{pmatrix}\in H.
\]
It represents a Coxeter element relative to the pinned torus $T_1$. Kostant's theorem makes it principal and regular. Hence $T_0=Z_G(n)^\circ$ and $T_1$ are in apposition, both $\Gamma$-stable.

There is also a stable Borel through $T_0$. The semisimple element $n\in H$ belongs to a torus $S\subset H$ \citep[Corollary~17.36, p.~363]{Milne2017}. Then $S\subset T_0$. No root of $(G,T_0)$ is trivial on $S$, since that would give value $1$ at $n$. Choose a real cocharacter vector of $S$ away from all root hyperplanes. It is $\Gamma$-fixed, and the roots positive on it define a $\Gamma$-stable Borel $B_0$. Together with the pinned $B_1$, this yields a fixed quadruple and a fixed generating principal element. Its full automorphism group is $\Gamma$ by the already proved inner rigidity.

For completeness, the construction also gives an effective test for every full automorphism group. Use this fixed quadruple as origin for the two Weyl sets. Then the class of $(w_0,w_1)$ has automorphism group
\[
 \{\gamma\in\Gamma:
 (\gamma(w_0)w_0^{-1},\gamma(w_1)w_1^{-1})
          \in A_h\times_{U_h}N_h\}.
\]
Indeed this membership says exactly that $\gamma(w_0,w_1)$ lies in its common-normalizer orbit; rigidity identifies that orbit stabilizer with the full automorphism group. In the excluded types $A_{2m}$ the diagram group has order two and no fixed class, so every full automorphism group is trivial.

\subsection{The type-\texorpdfstring{$A$}{A} normalizer in coordinates}
The Fourier model gives the converse and an explicit count of invariant classes. Put $I=\Z/n$. Every simultaneous normalizer is
\begin{equation}\label{eq:affine}
 [g_{a,b,d}],\qquad
 g_{a,b,d}e_j=\zeta^{ad j}e_{aj+b},
 \quad a\in I^\times,\quad b,d\in I.
\end{equation}
Its actions on coordinate and Fourier labels are respectively $j\mapsto aj+b$ and $c\mapsto a^{-1}c+d$.

To prove that the list is exhaustive, write an arbitrary common normalizer as
\[
 ge_j=u_je_{p(j)},\qquad gv_c=w_cv_{s(c)}.
\]
Comparison of coordinates, then division by the instances $c=0$ and $j=0$, gives
\[
 jc=(p(j)-p(0))(s(c)-s(0))\pmod n.
\]
At $j=c=1$ the two factors are units, even for composite $n$. Setting $a=p(1)-p(0)$ yields $p(j)=aj+b$ and $s(c)=a^{-1}c+d$; the remaining equation gives the scalar in \eqref{eq:affine}. Conversely direct substitution verifies the formula. Thus the common normalizer has $n^2\varphi(n)$ elements, as the general count predicts.

Inverse transpose $\vartheta(g)=(g^{-1})^{\mathsf T}$ preserves the torus pair. It sends the two orderings to
\[
 (\pi,\sigma)\longmapsto(\pi^{\mathrm{rev}},-\sigma^{\mathrm{rev}}),
\]
because $v_c^{\mathsf T}v_d=n\delta_{c+d,0}$. A class is fixed exactly when some $(a,b,d)$ satisfies
\begin{equation}\label{eq:reverse}
 a\pi(i)+b=\pi(n-1-i),\qquad
 a^{-1}\sigma(i)+d=-\sigma(n-1-i).
\end{equation}
For odd $n\geq3$, the first affine permutation must be an involution with exactly one fixed point. Thus $a^2=1$, and $a-1$ is a unit, forcing $a=-1$. The second condition then says a translation reverses an ordering. It is an involution, so $2d=0$ and therefore $d=0$, impossible for the identity permutation on at least three labels. This excludes $A_{2m}$. For even $n=2m$, an example satisfying \eqref{eq:reverse} is
\[
 \pi=(0,1,\ldots,m-1,2m-1,\ldots,m),\quad
 \sigma=(0,1,\ldots,2m-1),
\]
with $(a,b,d)=(1,m,1)$. Theorem~\ref{thm:main}(iii) follows.

\begin{proposition}\label{prop:fixedcount}
For $n=2m\geq4$, write $n=2^s u$ with $u$ odd, and let $\omega(u)$ count its distinct prime divisors. The number of outer-fixed classes in type $A_{n-1}$ is
\[
 F_n=\frac{(2^m m!)^2}{n^2\varphi(n)}
 \begin{cases}
 (n/2)2^{\omega(u)},&s=1,\\
 n2^{\omega(u)},&s\geq2.
 \end{cases}
\]
Their full automorphism groups are $C_2$, and all other classes have trivial automorphism group.
\end{proposition}
\begin{proof}
Since the normalizer acts freely, each fixed orbit contributes exactly $|K|$ pairs $(k,q)$ satisfying $kq=\vartheta q$. For $k=(a,b,d)$, this requires both
\[
 R(j)=aj+b,\qquad S(c)=-a^{-1}c-d
\]
to be conjugate to reversal, hence fixed-point-free involutions. Each has exactly $2^m m!$ compatible orderings. Thus $|K|F_n=t_n(2^m m!)^2$, where $t_n$ counts admissible triples.

The equations for involution are $a^2=1$, $(a+1)b=0$, $(a-1)d=0$ modulo $n$. At an odd prime power $p^e\Vert u$, the possibilities $a=\pm1$ each permit $p^e$ choices of $(b,d)$. Both resulting affine involutions have fixed points, as every involution of an odd set does. By the Chinese remainder theorem, fixed-point-freeness of both maps is therefore determined entirely modulo $2^s$; the odd part contributes $u2^{\omega(u)}$.

Modulo $2$, only $(a,b,d)=(1,1,1)$ works. Modulo $2^s$, $s\geq2$, $a=1$ requires $b=2^{s-1}$ and $d$ odd, giving $2^{s-1}$ choices; $a=-1$ gives another $2^{s-1}$. For $s\geq3$ the other square roots $1+2^{s-1}$ and $-1+2^{s-1}$ give none: in the first case $(a+1)b=0$ makes $b$ divisible by $2^{s-1}$, so $R$ has a fixed point; in the second, $(a-1)d=0$ similarly gives a fixed point of $S$. Multiplication of these local counts gives the formula. The automorphism assertion follows from rigidity.
\end{proof}

In particular $F_4=8$ out of $18$ and $F_6=192$ out of $7200$. The smallest obstruction is $A_2$: there are two classes, and inverse transpose exchanges them.

\section{What naturality does and does not select}\label{sec:natural}
The word ``distinguished'' in the letter does not specify a universal property. We make precise three possible requirements and determine their force.

\subsection{Isomorphisms and central isogenies}
For the groupoid of simple adjoint groups of one fixed type, $G\mapsto\mathcal X(G)$ is a functor to finite sets. A natural rule is a section of this functor. Evaluating at one group gives a fixed point of its automorphism group, whose action factors through $\Out(G)$. Conversely a fixed point $x\in\mathcal X(G)$ transports to any isomorphic $H$: define $x_H=f_*x$. Two choices of $f$ differ by an automorphism of $G$, so the definition is independent of the choice. This proves the first assertion of Theorem~\ref{thm:natural}.

For a central isogeny $\widetilde G\to G$ over $\C$, inverse image and image give inverse bijections on maximal tori and on Borel subgroups. The Lie algebra and root systems identify; thus apposition, defined through the adjoint principal element, and opposition are preserved. Inner-conjugacy classes also correspond, since the map on complex points is surjective. Therefore every choice on the adjoint group extends uniquely to its central forms. Compatibility with these central isogenies supplies no additional selection among classes.

\subsection{Opposition still leaves eight choices in \texorpdfstring{$A_3$}{A3}}
We strengthen the numerical example to show that the most immediate position requirement does not select a class.
\begin{proposition}\label{prop:A3}
For $\PGL_4$, exactly $16$ of the $18$ classes have opposite Borels, and all eight outer-fixed classes are among them.
\end{proposition}
\begin{proof}
Every entry of $F_4$ and $F_4^{-1}$ is nonzero, so only the middle $2\times2$ test in \eqref{eq:minor} matters. It fails exactly when the two complementary coordinate labels have the same parity and the first two Fourier labels have the same parity. Each ordering has this parity pattern in $2\cdot2!\cdot2!=8$ ways. Thus precisely $8^2=64$ Borel pairs fail. Dividing by $|K|=4^2\varphi(4)=32$ gives two nonopposite classes.

For a fixed class, the proof of Proposition~\ref{prop:fixedcount} shows that the affine reversal has either $a=1,b=2,d$ odd or $a=-1,d=2,b$ odd. In the first case $R$ is translation by $2$; the first half of any compatible coordinate ordering contains one member of each pair $\{0,2\},\{1,3\}$, so its complementary last two labels have opposite parity. In the second case $S$ is translation by $2$, so the first two Fourier labels have opposite parity. Either case makes the determinant nonzero. Hence all eight fixed classes are opposite.
\end{proof}

In particular maximal Bruhat length, equivalently the open opposite orbit, together with naturality under isomorphisms still leaves eight possibilities. A requirement based only on those two properties cannot be a uniqueness characterization.

\subsection{A field obstruction to split-real positivity}
\begin{proposition}\label{prop:field}
For a field $k\subset\C$, an apposition quadruple in the split group $\PGL_{n,k}$ with all four subgroup schemes defined over $k$ exists if and only if $k$ contains a primitive $n$th root of unity.
\end{proposition}
\begin{proof}
A $k$-Borel stabilizes a full $k$-flag. A maximal torus inside it is split: after lifting to upper triangular matrices in $\GL_n$, projection to the diagonal restricts to an injective homomorphism of $n$-dimensional tori, hence an isomorphism. Its kernel is trivial because it is both diagonalizable and unipotent. Both tori therefore admit eigenbases over $k$, with transition matrix $A\in\GL_n(k)$.

Over $\C$, the single apposition-pair orbit identifies $A$ with the Fourier matrix up to row and column permutations and nonzero scalings. Consequently all entries of $A$ are nonzero, and its cross-ratios
\[
 \frac{A_{ic}A_{jd}}{A_{id}A_{jc}}
\]
include a primitive $n$th root, since for the Fourier matrix rows $0,1$ and columns $0,1$ give $\zeta_n$. Scalings cancel and permutations only relabel the ratios. Hence $\zeta_n\in k$. Conversely, if $\zeta_n\in k$, the Fourier tori and their eigenline flags are all defined over $k$.
\end{proof}

Thus for $n\geq3$ no such quadruple is defined over the split real group. A proposal using totally positive split-real flags must first introduce different data; it cannot directly choose among these complex quadruples by placing all four subgroups in that real form. This obstruction does not exclude a positivity condition defined after some other construction.

The remaining geometric question is therefore precise: specify a property, beyond isomorphism equivariance and opposition, whose satisfying locus is a singleton in $\mathcal X(G)$ for each allowed type, and prove that it is preserved under the desired comparisons. Diagram folding is a possible additional comparison, but it is not an isomorphism or a central isogeny. Section~\ref{sec:folding} determines the standard symplectic comparison and shows that four choices remain already in type $A_3$. A unique selection compatible with all diagram foldings is not established here.

\section{The subregular predictions and the known rank theorem}\label{sec:subregular}
The second page of the letter predicts rank $r-1$ at subregular elements in the simply laced types, a smaller rank in the other types, and suggests that uniqueness of the subregular orbit might distinguish the first case. Both proposed distinctions fail.

\begin{theorem}[Slodowy; Richardson]\label{thm:subregular}
Let $\g$ be a complex simple Lie algebra of dimension $N$ and rank $r$. There is a unique nilpotent orbit of dimension $N-r-2$. At each of its points the differential of the adjoint quotient $\chi_{\g}:\g\to\tfr/W$ has rank $r-1$. The corresponding unipotent assertion holds for every connected complex group with this Lie algebra.
\end{theorem}

The nilpotent uniqueness is stated at the beginning of \citet[\S5, p.~234]{Richardson1987}; the rank is his Lemma~5.1 on the same page, explicitly attributed there to Slodowy's book \citep{Slodowy1980}. This is a known theorem, not an inference from the apposition classification. For $A_1$, the orbit is zero and the rank is zero.

For clarity, we give the passage to the group quotient, where the central form might otherwise cause confusion. The exponential and logarithm give inverse equivariant isomorphisms between the nilpotent and unipotent varieties. To see algebraicity, use their finite matrix series. If $u$ is unipotent in a faithful representation, the polynomial curve
\[
 u^t=\sum_{j\geq0}\binom tj(u-1)^j
\]
lies in the group at every nonnegative integer, hence everywhere; its tangent at zero is $\log u\in\g$. This also verifies the inverse direction without a choice of a matrix logarithm branch.

The ranks transfer as well. The torus exponential induces an analytic isomorphism $\psi$ between neighborhoods of $0$ in $\tfr/W$ and $1$ in $T/W$. By restriction of invariants, on a neighborhood of the nilpotent cone,
\begin{equation}\label{eq:ranktransfer}
 \chi_G(\exp x)=\psi(\chi_{\g}(x)).
\end{equation}
Indeed it holds on regular semisimple elements by conjugating to the torus, and extends by analytic density. The group restriction isomorphism holds without simple connectedness \citep[Corollary~6.4, p.~63]{Steinberg1965}. At nilpotent $e$, after left translation by $\exp(-e)$,
\[
 (d\exp)_e=\sum_{j\geq0}\frac{(-\operatorname{ad}e)^j}{(j+1)!},
\]
a finite polynomial with constant term $1$, hence invertible. Differentiating \eqref{eq:ranktransfer} proves equality of the group and Lie quotient ranks. Orbit dimensions are preserved by equivariance, giving the unipotent assertion in Theorem~\ref{thm:subregular}.

The smallest counterexample can be checked without a nilpotent-orbit table. In $\mathfrak{sp}_4$ put
\[
 J=\begin{pmatrix}0&I_2\\-I_2&0\end{pmatrix},\qquad
 e=\begin{pmatrix}0&I_2\\0&0\end{pmatrix}.
\]
For $Y=\left(\begin{smallmatrix}A&B\\C&-A^{\mathsf T}\end{smallmatrix}\right)$, with $B,C$ symmetric,
\[
 [Y,e]=\begin{pmatrix}-C&A+A^{\mathsf T}\\0&C\end{pmatrix}.
\]
The centralizer has dimension $3+1=4$, so the orbit has dimension $10-4=6=10-2-2$. The quotient coordinates are the characteristic coefficients $c_2=-\tr(Y^2)/2$ and $c_4=\det Y$. They generate the invariants because restriction to the diagonal torus gives the symmetric polynomials in the two squared eigenvalues, and regular semisimple elements are dense. At $e$,
\[
 (dc_2)_e(Y)=-\tr C\ne0,\qquad (dc_4)_e=0;
\]
the latter follows from the vanishing of every $3\times3$ minor of the rank-two matrix $e$. The quotient rank is $1=r-1$ despite the two root lengths. It is also $1$ at the image of $1+e$ in the adjoint group $\operatorname{PSp}_4$ by the rank transfer.

The separate questions on the same sheet about all quotient ranks, Springer-fibre dimensions, and rational singularities are broader than the subregular assertions treated here. None of the arguments above identifies the group quotient globally with the Lie-algebra quotient.

\section{Symplectic folding and the remaining ambiguity}\label{sec:folding}
We now determine exactly which classes arise under the standard symplectic folding, and show that this comparison still leaves more than one choice.

Folding and apposition already occur together in affine Toda theory; see \citet[Appendix~C.3--C.4, pp.~62--65]{DoreyPolvara2022} for Cartans and twisted Coxeter elements. Here the comparison concerns the finite sets of classes equipped with two Borel subgroups.

\begin{proposition}\label{prop:folding}
Let $n=2m\geq4$, $G=\PGL_n(\C)$, and let $H=\operatorname{PSp}_n(\C)$ be the fixed group of a symplectic outer involution $\sigma$ of $G$. For either group use $\mathcal X$ for the inner classes of apposition quadruples. Let $\mathcal Y\subset\mathcal X(G)$ consist of the outer-fixed classes whose unique nontrivial outer automorphism fixes the intrinsic cyclic group $T_0\cap N_G(T_1)$ pointwise. Extending the tori to their centralizers in $G$ and extending their isotropic flags to complete selfdual flags induces a bijection
\[
 \mathcal X(H)\longrightarrow\mathcal Y.
\]
In particular
\[
 |\mathcal Y|=\frac{(2^m m!)^2}{2n\varphi(n)}.
\]
For $n=4$ this standard symplectic folding condition leaves four classes.
\end{proposition}
\begin{proof}
A maximal torus of $H$ has the distinct standard weights $\varepsilon_1,\ldots,\varepsilon_m,-\varepsilon_1,\ldots,-\varepsilon_m$, so its centralizer in $G$ is a maximal torus. An isotropic complete flag in $\C^n$ extends uniquely to a complete selfdual flag. Hence the extension of a quadruple is well defined and equivariant under $H$.

Conjugate its apposition pair in $H$ to the cyclic-shift model of Section~\ref{sec:classical}. The extended tori are the full diagonal and Fourier tori in $G$, so they are in apposition there. The diagonal principal element belongs to the projective symplectic group: a scalar multiple of $\diag(1,\zeta,\ldots,\zeta^{n-1})$ preserves the alternating form. Thus $\sigma$ fixes the extended quadruple and its generating principal element. The map lands in $\mathcal Y$.

If $g\in G$ conjugates two extended quadruples, then $g\sigma g^{-1}$ and $\sigma$ are automorphisms of the second quadruple with the same outer class. Their quotient is an inner stabilizer, hence is trivial. Therefore $g$ is fixed by $\sigma$. In $\PGL_n$, that fixed group is $H$: a projective symplectic similitude becomes symplectic after multiplying a lift by a square root of its multiplier. Thus $g\in H$, proving injectivity.

To prove surjectivity it suffices to count. In the Fourier model, if $g_{a,b,d}$ carries a quadruple to its inverse-transpose image, its unique outer automorphism is $\operatorname{Int}(g_{a,b,d}^{-1})\vartheta$. On the principal diagonal element it acts by exponent $-a$. Pointwise fixation of the cyclic group is therefore exactly $a=-1$ in $\Z/n$. The two affine reversals in Proposition~\ref{prop:fixedcount} then become $j\mapsto-j+b$ and $c\mapsto c-d$. They are fixed-point-free involutions exactly when $b$ is odd and $d=n/2$. There are $n/2$ such triples, each admitting $(2^m m!)^2$ ordering pairs. Freeness gives
\[
 |\mathcal Y|=\frac{(n/2)(2^m m!)^2}{n^2\varphi(n)}.
\]
On the other hand $H$ has Weyl-group order $2^m m!$, Coxeter number $n$, and simply connected center of order two. Theorem~\ref{thm:main}(i) gives exactly the same cardinality for $\mathcal X(H)$. The injective map is consequently onto.
\end{proof}

This shows that compatibility with this particular folding does not by itself distinguish a class: in rank two the folded group $\operatorname{PSp}_4$ already has four choices and has no outer automorphisms. The question of a geometric singleton condition compatible with all desired foldings remains open.

\section*{Method and records}
The mathematics and text were produced by GPT-6 Astra (ChatGPT Pro mode). The main counting, opposition and outer-symmetry statements, together with the split type-$A$ field criterion, were independently reconstructed from their statements by a separate instance of the same model before it read the proposed proofs. That reconstruction retained Kostant's cyclic-element theorem and standard Lie theory as explicit inputs. Another same-model statement-only derivation proved the symplectic-folding bijection and count and the eight outer-fixed opposite classes in type $A_3$. All written proofs also underwent targeted checks, including the exceptional integer reflection calculations and an exhaustive $A_3$ flag enumeration. Recorded checks by Claude Opus~5.5 at effort high include the low-rank counts and the $\mathfrak{sp}_4$ counterexample. These are machine checks and mathematical reviews, not a human verification. No human has checked this paper.

The text is dedicated to the public domain under CC0~1.0. The program \emph{Apposition Certificates} and mathematical check records are available through the workbench, \url{\workbenchurl}. The permanent address for the paper's records is \url{\recordurl}.

\begin{thebibliography}{99}
\bibitem[Dorey and Polvara(2022)]{DoreyPolvara2022}
Dorey, P., \& Polvara, D. (2022). Tree level integrability in 2d quantum field theories and affine Toda models.
\emph{Journal of High Energy Physics, 2022}(2), Article~199.
\url{https://doi.org/10.1007/JHEP02(2022)199}.
Author version (the page locators used here): \url{https://arxiv.org/abs/2111.02210}.

\bibitem[Elashvili et~al.(2013)]{Elashvili2013}
Elashvili, A.~G., Kac, V.~G., \& Vinberg, E.~B. (2013).
Cyclic elements in semisimple Lie algebras.
\emph{Transformation Groups, 18}, 97--130.
\url{https://arxiv.org/abs/1205.0515}.

\bibitem[Etingof(n.d.)]{Etingof}
Etingof, P. (n.d.). \emph{18.745: Lie groups and Lie algebras, I} [Course notes]. Massachusetts Institute of Technology.
\url{https://math.mit.edu/~etingof/lnlg.pdf}.

\bibitem[Grothendieck(1969)]{Grothendieck1969}
Grothendieck, A. (1969, October 22). \emph{Letter to B. Kostant} [Manuscript]. Universit\'e de Montpellier, cote~123, scans~22--23, folios~21--22.
\url{https://grothendieck.umontpellier.fr/123.pdf}.

\bibitem[Guest and Ho(2019)]{GuestHo2019}
Guest, M.~A., \& Ho, N.-K. (2019). Kostant, Steinberg, and the Stokes matrices of the $tt^*$-Toda equations.
\emph{Selecta Mathematica, 25}, Article~50.
\url{https://doi.org/10.1007/s00029-019-0494-7}.
Author version (the page locators used here): \url{https://arxiv.org/abs/1802.01126}.

\bibitem[Hua(2026)]{Hua2026}
Hua, M. (2026). \emph{\'Enonc\'es des notes de travail de Grothendieck d\'emontr\'es en Lean~4, et questions ouvertes} (Version~2) [Preprint]. Hexagon.
\url{https://hexagonmath.org/2611.00125v2}.

\bibitem[Kostant(1959)]{Kostant1959}
Kostant, B. (1959). The principal three-dimensional subgroup and the Betti numbers of a complex simple Lie group.
\emph{American Journal of Mathematics, 81}, 973--1032.
\url{https://doi.org/10.2307/2372999}.

\bibitem[Kostant(2010)]{Kostant2010}
Kostant, B. (2010). Experimental evidence for the occurrence of $E_8$ in nature and the radii of the Gosset circles.
\emph{Selecta Mathematica, 16}, 419--438.
\url{https://doi.org/10.1007/s00029-010-0023-1}.
Author version (the page locators used here):
\url{https://arxiv.org/abs/1003.0046v1}.

\bibitem[Milne(2017)]{Milne2017}
Milne, J.~S. (2017). \emph{Algebraic groups: The theory of group schemes of finite type over a field}. Cambridge University Press.
\url{https://www.jmilne.org/math/Books/iAG2017.pdf}.

\bibitem[Richardson(1987)]{Richardson1987}
Richardson, R.~W. (1987). Derivatives of invariant polynomials on a semisimple Lie algebra.
In \emph{Miniconference on harmonic analysis and operator algebras} (Proceedings of the Centre for Mathematical Analysis, Australian National University, Vol.~15, pp.~228--241).
\url{https://maths.anu.edu.au/files/CMAProcVol15-Richardson_1.pdf}.

\bibitem[Slodowy(1980)]{Slodowy1980}
Slodowy, P. (1980). \emph{Simple singularities and simple algebraic groups} (Lecture Notes in Mathematics, Vol.~815). Springer.
\url{https://doi.org/10.1007/BFb0090294}.

\bibitem[Steinberg(1965)]{Steinberg1965}
Steinberg, R. (1965). Regular elements of semisimple algebraic groups.
\emph{Publications Math\'ematiques de l'IH\'ES, 25}, 49--80.
\url{https://www.numdam.org/item/PMIHES_1965__25__49_0/}.

\bibitem[Tao(2005)]{Tao2005}
Tao, T. (2005). An uncertainty principle for cyclic groups of prime order.
\emph{Mathematical Research Letters, 12}(1), 121--127.
\url{https://arxiv.org/abs/math/0308286}.
\end{thebibliography}
\end{document}
